% =========================================================
% REGLAS DE NEGOCIO - SISTEMA DE GESTIÓN HOTEL PARA PERROS
% =========================================================

\begin{BussinesRule}{BR001}{Fecha de Nacimiento válida para Mascota}
	\BRitem[Tipo:] Regla de integridad referencial o estructural.
	\BRitem[Clase:] Habilitadora.
	\BRitem[Nivel:] Control.
	\BRitem[Descripción:] La fecha de nacimiento de una mascota debe ser posterior al 1 de Enero de 1990 y anterior a la fecha actual. Esto asegura que se registren mascotas reales y evita errores en el cálculo de edad.
	\BRitem[Motivación:] Prevenir errores en el registro de mascotas con fechas ilógicas y garantizar la integridad de los datos de edad para cálculos de servicios veterinarios y tarifas.
	\BRitem[Sentencia:] $\forall m \in Mascota \Rightarrow 01\textrm{-}Enero\textrm{-}1990~<~m.fechaNacimiento~<~fechaActual$.
	\BRitem[Ejemplo positivo:] Para el día 11 de Enero del 2026, cumplen la regla:
		\begin{itemize}
			\item 10 de Enero del 2026
			\item 15 de Marzo del 2020
			\item 5 de Mayo del 2000
			\item 2 de Enero del 1990
		\end{itemize}
	\BRitem[Ejemplo negativo:] Para el día 11 de Enero del 2026, NO cumplen:
		\begin{itemize}
			\item 11 de Enero del 2026 (fecha actual no permitida)
			\item 20 de Diciembre del 2026 (fecha futura)
			\item 1 de Enero del 1990 (límite inferior no incluido)
			\item 31 de Diciembre del 1989
		\end{itemize}
	\BRitem[Referenciado por:] CUCE-PET-001, CUCE-PET-002.
\end{BussinesRule}

\begin{BussinesRule}{BR002}{Validación de fechas de Reservación}
	\BRitem[Tipo:] Regla de integridad referencial o estructural.
	\BRitem[Clase:] Habilitadora.
	\BRitem[Nivel:] Control.
	\BRitem[Descripción:] En toda reservación, la fecha de fin debe ser estrictamente posterior a la fecha de inicio. Además, la fecha de inicio debe ser igual o posterior a la fecha actual.
	\BRitem[Motivación:] Garantizar la lógica temporal de las reservaciones y evitar conflictos en la asignación de habitaciones.
	\BRitem[Sentencia:] $\forall r \in Reservacion \Rightarrow fechaActual \leq r.fechaInicio~<~r.fechaFin$.
	\BRitem[Ejemplo positivo:] Para una reservación creada el 11 de Enero del 2026:
		\begin{itemize}
			\item Inicio: 11/01/2026, Fin: 15/01/2026
			\item Inicio: 20/01/2026, Fin: 25/01/2026
			\item Inicio: 11/01/2026, Fin: 12/01/2026 (una noche)
		\end{itemize}
	\BRitem[Ejemplo negativo:] Para una reservación creada el 11 de Enero del 2026:
		\begin{itemize}
			\item Inicio: 15/01/2026, Fin: 10/01/2026 (fin antes de inicio)
			\item Inicio: 15/01/2026, Fin: 15/01/2026 (misma fecha)
			\item Inicio: 05/01/2026, Fin: 10/01/2026 (inicio en el pasado)
		\end{itemize}
	\BRitem[Referenciado por:] CUCE-RES-001, CUCE-RES-002.
\end{BussinesRule}

\begin{BussinesRule}{BR003}{Disponibilidad de Habitación}
	\BRitem[Tipo:] Regla de inferencia de un hecho.
	\BRitem[Clase:] Habilitadora.
	\BRitem[Nivel:] Control.
	\BRitem[Descripción:] Una habitación solo puede ser asignada a una reservación si no existe otra reservación confirmada que tenga traslape de fechas con el periodo solicitado. Se considera traslape cuando hay intersección entre $[fechaInicio_1, fechaFin_1]$ y $[fechaInicio_2, fechaFin_2]$.
	\BRitem[Motivación:] Evitar doble asignación de habitaciones y garantizar la disponibilidad del espacio para cada mascota.
	\BRitem[Sentencia:] Sea $R$ el conjunto de reservaciones confirmadas en habitación $h$:
		\begin{displaymath}
			Disponible(h, f_i, f_f) = \neg\exists r \in R : (r.fechaInicio < f_f) \land (r.fechaFin > f_i)
		\end{displaymath}
	\BRitem[Ejemplo positivo:] Habitación disponible cuando:
		\begin{itemize}
			\item Nueva reserva: 15-20 Ene, Reservas existentes: 10-14 Ene, 21-25 Ene
			\item Nueva reserva: 10-15 Mar, No hay reservas en ese periodo
		\end{itemize}
	\BRitem[Ejemplo negativo:] Habitación NO disponible cuando:
		\begin{itemize}
			\item Nueva: 15-20 Ene, Existe: 18-22 Ene (traslape parcial)
			\item Nueva: 10-25 Ene, Existe: 15-20 Ene (traslape completo)
			\item Nueva: 15-20 Ene, Existe: 15-20 Ene (mismas fechas)
		\end{itemize}
	\BRitem[Referenciado por:] CUCE-RES-003, CUCE-HAB-001.
\end{BussinesRule}

\begin{BussinesRule}{BR004}{Autorización de Reservación por Propietario}
	\BRitem[Tipo:] Regla de integridad referencial o estructural.
	\BRitem[Clase:] Habilitadora.
	\BRitem[Nivel:] Control.
	\BRitem[Descripción:] Un propietario solo puede crear o modificar reservaciones para mascotas que le pertenecen. Los administradores pueden gestionar reservaciones de cualquier mascota.
	\BRitem[Motivación:] Garantizar la privacidad y seguridad de la información, evitando accesos no autorizados a reservaciones de terceros.
	\BRitem[Sentencia:] 
		\begin{displaymath}
			Autorizado(p, r) = (p.propietarioId = r.mascota.propietarioId) \lor (p.rol = \textrm{Admin})
		\end{displaymath}
	\BRitem[Ejemplo positivo:]
		\begin{itemize}
			\item Propietario ID 5 crea reservación para su mascota ID 12
			\item Admin crea reservación para mascota ID 20 de cualquier propietario
			\item Propietario ID 3 edita su propia reservación ID 45
		\end{itemize}
	\BRitem[Ejemplo negativo:]
		\begin{itemize}
			\item Propietario ID 5 intenta crear reservación para mascota ID 18 (de otro propietario)
			\item Usuario no-admin intenta editar reservación de otro propietario
		\end{itemize}
	\BRitem[Referenciado por:] CUCE-RES-001, CUCE-RES-002, CUCE-RES-004.
\end{BussinesRule}

\begin{BussinesRule}{BR005}{Peso mínimo y máximo de Mascota}
	\BRitem[Tipo:] Regla de integridad referencial o estructural.
	\BRitem[Clase:] Habilitadora.
	\BRitem[Nivel:] Control.
	\BRitem[Descripción:] El peso de una mascota debe estar en el rango de 0.1 kg a 200 kg. Este rango cubre desde mascotas pequeñas (como chihuahuas) hasta razas muy grandes.
	\BRitem[Motivación:] Evitar errores de captura y validar que los datos ingresados sean realistas para calcular tarifas y requerimientos de espacio.
	\BRitem[Sentencia:] $\forall m \in Mascota \Rightarrow 0.1 \leq m.pesoKg \leq 200$.
	\BRitem[Ejemplo positivo:]
		\begin{itemize}
			\item Chihuahua: 2.5 kg
			\item Labrador: 30 kg
			\item Gran Danés: 80 kg
			\item San Bernardo: 150 kg
		\end{itemize}
	\BRitem[Ejemplo negativo:]
		\begin{itemize}
			\item 0 kg (peso nulo)
			\item 0.05 kg (por debajo del mínimo)
			\item 250 kg (excede límite superior)
			\item -5 kg (valor negativo)
		\end{itemize}
	\BRitem[Referenciado por:] CUCE-PET-001, CUCE-PET-003.
\end{BussinesRule}

\begin{BussinesRule}{BR006}{Vigencia de Vacunación}
	\BRitem[Tipo:] Regla de operación (calcular o determinar un valor).
	\BRitem[Clase:] Habilitadora.
	\BRitem[Nivel:] Influencia.
	\BRitem[Descripción:] Una vacuna se considera vigente si su fecha de vigencia es posterior o igual a la fecha actual. Las vacunas con vigencia vencida deben ser señaladas en el historial médico.
	\BRitem[Motivación:] Alertar a los propietarios sobre vacunas próximas a vencer o vencidas para mantener la salud de las mascotas y cumplir requisitos de hospedaje.
	\BRitem[Sentencia:] 
		\begin{displaymath}
			VacunaVigente(v) = \left\{ \begin{array}{ll}
			\textrm{Vigente} & , \textrm{si}~v.vigenciaHasta \geq fechaActual\\
			\textrm{Vencida} & , \textrm{si}~v.vigenciaHasta < fechaActual
			\end{array} \right.
		\end{displaymath}
	\BRitem[Ejemplo positivo:] Para la fecha actual 11/01/2026:
		\begin{itemize}
			\item Vacuna antirrábica vigente hasta 15/06/2026 (Vigente)
			\item Vacuna polivalente vigente hasta 11/01/2026 (Vigente, límite)
		\end{itemize}
	\BRitem[Ejemplo negativo:] Para la fecha actual 11/01/2026:
		\begin{itemize}
			\item Vacuna antirrábica vigente hasta 10/01/2026 (Vencida)
			\item Vacuna polivalente vigente hasta 15/12/2025 (Vencida)
		\end{itemize}
	\BRitem[Referenciado por:] CUCE-VAC-001, CUCE-PET-004.
\end{BussinesRule}

\begin{BussinesRule}{BR007}{Cálculo de Monto Total de Reservación}
	\BRitem[Tipo:] Regla de operación (calcular o determinar un valor).
	\BRitem[Clase:] Habilitadora.
	\BRitem[Nivel:] Control.
	\BRitem[Descripción:] El monto total de una reservación se calcula como la suma del costo de hospedaje más el costo de todos los servicios adicionales contratados durante la estancia.
	\BRitem[Motivación:] Proporcionar transparencia en los costos y facilitar la facturación al cliente.
	\BRitem[Sentencia:] Sea $S$ el conjunto de servicios en reservación $r$:
		\begin{displaymath}
			MontoTotal(r) = r.montoHospedaje + \sum_{s \in S} (s.precioAlMomento \times s.cantidad)
		\end{displaymath}
	\BRitem[Ejemplo positivo:]
		\begin{itemize}
			\item Hospedaje: \$500, Baño: \$150 $\times$ 1, Paseo: \$80 $\times$ 3 = \$890
			\item Hospedaje: \$1200, Sin servicios = \$1200
		\end{itemize}
	\BRitem[Ejemplo negativo:] Valores que generarían error:
		\begin{itemize}
			\item Monto negativo de hospedaje o servicios
			\item Cantidad de servicio igual a cero
			\item Sumar servicios de otra reservación
		\end{itemize}
	\BRitem[Referenciado por:] CUCE-RES-005, CUCE-PAG-001.
\end{BussinesRule}

\begin{BussinesRule}{BR008}{Validación de horario de Citas de Servicio}
	\BRitem[Tipo:] Regla de operación (calcular o determinar un valor).
	\BRitem[Clase:] Habilitadora.
	\BRitem[Nivel:] Control.
	\BRitem[Descripción:] Las citas de servicio (veterinaria, peluquería, entrenamiento) asignadas a un mismo empleado no pueden traslaparse en horario. Dos citas se traslapan si el inicio de una es antes del fin de la otra y viceversa.
	\BRitem[Motivación:] Optimizar la agenda de empleados y evitar conflictos de disponibilidad para brindar servicios de calidad.
	\BRitem[Sentencia:] Sea $C$ el conjunto de citas del empleado $e$:
		\begin{displaymath}
			\begin{array}{l}
			NoTraslape(c_1, c_2) = \\
			\quad (c_1.fechaHoraFin \leq c_2.fechaHoraInicio) \lor (c_2.fechaHoraFin \leq c_1.fechaHoraInicio)
			\end{array}
		\end{displaymath}
	\BRitem[Ejemplo positivo:] Citas que NO traslapan:
		\begin{itemize}
			\item Cita 1: 10:00-11:30, Cita 2: 12:00-13:30
			\item Cita 1: 09:00-10:00, Cita 2: 10:00-11:00 (consecutivas)
		\end{itemize}
	\BRitem[Ejemplo negativo:] Citas que SÍ traslapan:
		\begin{itemize}
			\item Cita 1: 10:00-12:00, Cita 2: 11:00-13:00
			\item Cita 1: 09:00-11:00, Cita 2: 10:30-12:00
		\end{itemize}
	\BRitem[Referenciado por:] CUCE-CITA-001, CUCE-EMP-002.
\end{BussinesRule}

\begin{BussinesRule}{BR009}{Cálculo de días de Hospedaje}
	\BRitem[Tipo:] Regla de operación (calcular o determinar un valor).
	\BRitem[Clase:] Habilitadora.
	\BRitem[Nivel:] Control.
	\BRitem[Descripción:] El número de días de hospedaje se calcula como la diferencia en días completos entre la fecha de fin y la fecha de inicio de la reservación. Si la estancia es de menos de 24 horas, se cuenta como un día completo.
	\BRitem[Motivación:] Estandarizar el cálculo de días para la facturación y evitar ambigüedades en cobros.
	\BRitem[Sentencia:] 
		\begin{displaymath}
			DiasHospedaje(r) = \max(1, \lceil (r.fechaFin - r.fechaInicio) / 86400 \rceil)
		\end{displaymath}
	\BRitem[Ejemplo positivo:]
		\begin{itemize}
			\item Check-in: 15/01 10:00, Check-out: 20/01 10:00 = 5 días
			\item Check-in: 15/01 14:00, Check-out: 15/01 20:00 = 1 día
			\item Check-in: 10/02 09:00, Check-out: 12/02 15:00 = 3 días
		\end{itemize}
	\BRitem[Ejemplo negativo:] No aplican ejemplos negativos, pero resultados incorrectos serían:
		\begin{itemize}
			\item Calcular 0 días para estancia de horas
			\item No redondear al alza días parciales
		\end{itemize}
	\BRitem[Referenciado por:] CUCE-RES-006, CUCE-PAG-002.
\end{BussinesRule}

\begin{BussinesRule}{BR010}{Estados válidos de Reservación}
	\BRitem[Tipo:] Regla de integridad referencial o estructural.
	\BRitem[Clase:] Habilitadora.
	\BRitem[Nivel:] Control.
	\BRitem[Descripción:] Toda reservación debe tener un estado válido del catálogo de estados. Los estados permitidos son: Pendiente, Confirmada, En curso, Completada, Cancelada. Las transiciones entre estados deben seguir un flujo lógico.
	\BRitem[Motivación:] Mantener control del ciclo de vida de las reservaciones y facilitar reportes y estadísticas.
	\BRitem[Sentencia:] 
		\begin{displaymath}
			\forall r \in Reservacion \Rightarrow r.estadoId \in \{1, 2, 3, 4, 5\}
		\end{displaymath}
		Transiciones válidas: Pendiente $\rightarrow$ Confirmada/Cancelada, Confirmada $\rightarrow$ En~curso/Cancelada, En~curso $\rightarrow$ Completada
	\BRitem[Ejemplo positivo:]
		\begin{itemize}
			\item Nueva reserva creada en estado "Pendiente"
			\item Reserva pasa de "Confirmada" a "En curso" en check-in
			\item Reserva pasa de "En curso" a "Completada" en check-out
		\end{itemize}
	\BRitem[Ejemplo negativo:]
		\begin{itemize}
			\item Crear reserva con estado "Completada"
			\item Pasar de "Pendiente" directamente a "En curso"
			\item Usar estado inexistente (ID 99)
		\end{itemize}
	\BRitem[Referenciado por:] CUCE-RES-007, CUCE-RES-008.
\end{BussinesRule}
